% =====================================================================
% master/preamble.tex
% Shared layout and machinery for resume.master.tex and
% cover-letter.master.tex. Loaded right after \documentclass.
% You shouldn't need to edit this file to apply for a job.
% =====================================================================

% ---------------------------------------------------------------------
% 1. Page and text layout (tested fixes from the Résumé Audit, 2026-09-11)
% ---------------------------------------------------------------------
\geometry{letterpaper, left=1.5cm, right=1.5cm, top=1.1cm, bottom=1.2cm, footskip=.5cm}
\fontdir[fonts/]
\definecolor{awesome}{HTML}{222A68}
\definecolor{flagcolor}{HTML}{D35400}
\setbool{acvSectionColorHighlight}{false}
\XeTeXinterwordspaceshaping=2                 % real spaces in the PDF text (pypdf reads words, not "Datascientistwith...")
%% ...but shaping=2 also merges consecutive \href annotations on one line: the
%% second link's start is swallowed, so the first link's hotspot stretches over
%% the whole line. Every header link pointed at LinkedIn because of this, and
%% tel:/mailto: produced no annotation at all. An \mbox around each link keeps
%% its annotation intact. Use \slink, never a bare \href, wherever two or more
%% links share a line.
\newcommand{\slink}[2]{\mbox{\href{#1}{#2}}}
\hyphenpenalty=10000 \exhyphenpenalty=10000   % never split a word across lines
\setlength{\parskip}{0pt}
\setlength{\parindent}{0pt}
\pagestyle{empty}                              % no footer: parsers read footers as body text
\raggedright
\usepackage{needspace}                        % keep a role's header with its first bullet

% ---------------------------------------------------------------------
% 2. Build modes
%    Catalog  : \archetype = all  -> every bullet, with its ID and tags
%    Base     : \archetype = ds | mle | de | rs | uxr
%    Final    : application builds (\FinalBuildFlag defined) -> no highlights,
%               but every unverified fact / unfilled slot is logged and the
%               build script refuses to finish.
% ---------------------------------------------------------------------
\providecommand{\archetype}{all}
\newif\ifcatalog
\ifdefstring{\archetype}{all}{\catalogtrue}{\catalogfalse}
\newif\ifFinalBuild
\ifdefined\FinalBuildFlag \FinalBuildtrue \else \FinalBuildfalse \fi

% ---------------------------------------------------------------------
% One-page mode: \def\OnePageFlag{1} in an application file (or
% scripts/build.sh bases --onepage). Prints the same facts, less of them:
% each role keeps only its first \onepagemax bullets, roles wrapped in
% \unlessonepage are dropped, and publications shrink to the summary line.
% A bullet an application force-includes is never dropped.
% ---------------------------------------------------------------------
\newif\ifonepage
\ifdefined\OnePageFlag \onepagetrue \else \onepagefalse \fi
\newcounter{bulincluded}
\def\OnePageMax{2}
\newcommand{\onepagemax}[1]{\def\OnePageMax{#1}}
\newcommand{\unlessonepage}[1]{\ifonepage\else#1\fi}

% \unverified[label]{value}: a fact that disagrees between old files or needs your records
\newcommand{\unverified}[2][]{%
  \ifFinalBuild #2\else\textcolor{flagcolor}{#2\textsuperscript{?}}\fi
  \typeout{UNVERIFIED FACT: #1 = \detokenize{#2}}}
% \slot{what to write}: a placeholder that must be filled before sending
\newcommand{\slot}[1]{%
  \textcolor{red}{\textbf{[[SLOT: #1]]}}%
  \typeout{UNFILLED SLOT: \detokenize{#1}}}

% ---------------------------------------------------------------------
% 3. Role-type (archetype) helpers
% ---------------------------------------------------------------------
\newcommand{\archcheck}[1]{\ifdefstring{\archetype}{#1}{\def\archhit{1}}{}}
% \ifarch{ds,mle}{if the current role type is listed}{otherwise}
\newcommand{\ifarch}[3]{\def\archhit{0}\forcsvlist{\archcheck}{#1}\ifdefstring{\archhit}{1}{#2}{#3}}
% \onlyfor{rs,uxr}{content}: printed for those role types and in the catalog
\newcommand{\onlyfor}[2]{\ifcatalog #2\else\ifarch{#1}{#2}{}\fi}
% \catalogonly{content}: printed only in the catalog
\newcommand{\catalogonly}[1]{\ifcatalog #1\fi}

% Per-application include/exclude lists (set in an application file)
\providecommand{\ExcludeIDs}{}
\providecommand{\IncludeIDs}{}
\def\excludedidslist{}
\def\includedidslist{}
\renewcommand*{\do}[1]{\listadd\excludedidslist{#1}}
\expandafter\docsvlist\expandafter{\ExcludeIDs}
\renewcommand*{\do}[1]{\listadd\includedidslist{#1}}
\expandafter\docsvlist\expandafter{\IncludeIDs}

% ---------------------------------------------------------------------
% 4. Bullets: \bul{id}{role types}{text} ... \printbullets
%    A bullet prints when the current role type is in its list (or in the
%    catalog), unless the application excludes its ID; an application can
%    force-include any ID. An application can also reword one bullet with
%    \expandafter\def\csname override@<id>\endcsname{new text}.
% ---------------------------------------------------------------------
\newcommand{\bulletbuffer}{}
\newcommand{\bulcheckfor}[1]{\ifdefstring{\archetype}{#1}{\def\bulhit{1}}{}}
\newcommand{\catalogtag}[2]{\ifcatalog\hspace{0.35em}{\fontsize{7.5pt}{1em}\bodyfont\color{gray}[#1\,\textbar\,#2]}\fi}
\newcommand{\bul}[3]{%
  \def\bulhit{0}%
  \ifcatalog\def\bulhit{1}\else\forcsvlist{\bulcheckfor}{#2}\fi
  \ifinlist{#1}{\excludedidslist}{\def\bulhit{0}}{}%
  \ifonepage
    \ifdefstring{\bulhit}{1}%
      {\ifnum\value{bulincluded}<\OnePageMax\relax\else\def\bulhit{0}\fi}{}%
  \fi
  \ifinlist{#1}{\includedidslist}{\def\bulhit{1}}{}%
  \ifdefstring{\bulhit}{1}{%
    \stepcounter{bulincluded}%
    \ifcsdef{override@#1}%
      {\appto\bulletbuffer{\item \csuse{override@#1}\catalogtag{#1}{#2}}}%
      {\appto\bulletbuffer{\item #3\catalogtag{#1}{#2}}}}{}%
}
\newenvironment{rolebullets}{%
  \begin{itemize}[leftmargin=1.05em, labelsep=0.45em, itemsep=1.2pt, parsep=0pt, topsep=1.5pt, partopsep=0pt, label={\textbullet}]%
  \raggedright\fontsize{10pt}{12.4pt}\bodyfont\color{text}}%
  {\end{itemize}}
\newcommand{\printbullets}{%
  \ifdefempty{\bulletbuffer}{}{\begin{rolebullets}\bulletbuffer\end{rolebullets}}%
  \gdef\bulletbuffer{}\setcounter{bulincluded}{0}\def\OnePageMax{2}}

% ---------------------------------------------------------------------
% 5. Layout macros: header, sections, roles, skills, publications
% ---------------------------------------------------------------------
\def\theheadline{}
\def\thesummary{}
\def\theskills{}
\newcommand{\catalogentry}[2]{\par\noindent{\fontsize{8.5pt}{1em}\bodyfont\bfseries\color{awesome}[#1]}\hspace{0.4em}#2\par\vspace{3pt}}
% Headline (line under your name), summary and skills: one per role type
\newcommand{\headlinefor}[2]{\ifarch{#1}{\gdef\theheadline{#2}}{}}
\newcommand{\summaryfor}[2]{\ifcatalog\catalogentry{#1}{#2}\else\ifarch{#1}{\gdef\thesummary{#2}}{}\fi}
\newcommand{\skillsfor}[2]{\ifcatalog\catalogentry{#1}{}#2\vspace{2pt}\else\ifarch{#1}{\gdef\theskills{#2}}{}\fi}
\newcommand{\printheadline}{%
  \ifdefined\HeadlineOverride\HeadlineOverride
  \else\ifcatalog\ifdefined\CatalogHeadline\CatalogHeadline\else MASTER CATALOG \textbar{} every bullet for every role type \textbar{} not for sending\fi%
  \else\theheadline\fi\fi}
\newcommand{\printsummary}{%
  \ifcatalog\else{\fontsize{10pt}{12.6pt}\bodyfont\color{text}%
    \ifdefined\SummaryOverride\SummaryOverride\else\thesummary\fi\par}\fi}
\newcommand{\printskills}{\ifcatalog\else\ifdefined\SkillsOverride\SkillsOverride\else\theskills\fi\fi}

\newcommand{\masterheader}{%
  \begin{center}
    {\fontsize{25pt}{25pt}\headerfont\bfseries\color{darktext}\FactName\FactPostNominal}\par\vspace{2.2mm}
    {\fontsize{10.5pt}{13pt}\bodyfont\color{awesome}\printheadline}\par\vspace{1.4mm}
    {\fontsize{9.5pt}{12pt}\bodyfont\color{text}\FactCity\ifdefined\FactPhone\ \textbar\ \FactPhone\fi\ \textbar\ \FactEmail}\par\vspace{0.5mm}
    {\fontsize{9.5pt}{12pt}\bodyfont\color{text}\FactLinkedIn\ \textbar\ \FactGitHub\ \textbar\ \FactWebsite\onlyfor{rs,uxr}{\ \textbar\ \FactScholar}}\par
  \end{center}\vspace{-1mm}}

\newcommand{\msection}[1]{%
  \par\ifonepage\vspace{2pt}\else\vspace{6pt}\fi
  \noindent{\fontsize{12.5pt}{14pt}\bodyfont\bfseries\color{awesome}#1}\par\ifonepage\vspace{1pt}\else\vspace{1.5pt}\fi
  \noindent{\color{lightgray}\rule{\linewidth}{0.6pt}}\par\ifonepage\vspace{1.5pt}\else\vspace{2pt}\fi}

% \cvrole{employer}{location}{title}{dates}  -- then \bul ... \printbullets
\newcommand{\cvrole}[4]{%
  \par\needspace{5\baselineskip}\ifonepage\vspace{2.5pt}\else\vspace{4.5pt}\fi
  \noindent{\fontsize{10.5pt}{13pt}\bodyfont\bfseries\color{darktext}#1}\hfill{\fontsize{9.5pt}{13pt}\bodyfont\color{darktext}#2}\par\nopagebreak
  \noindent{\fontsize{10pt}{12.5pt}\bodyfont\color{darktext}#3}\hfill{\fontsize{9.5pt}{12.5pt}\bodyfont\color{darktext}#4}\par\nopagebreak}

% A second (earlier) appointment with the same employer: prints title + dates
% under the employer already printed by \cvrole. Harvard Law's resume guidance:
% "If you have held more than one job with the same employer, enter both jobs
% under one header." Keeps one employer/one location for a parser to read.
\newcommand{\cvsubrole}[2]{%
  \par\needspace{3\baselineskip}%
  \noindent{\fontsize{10pt}{12.5pt}\bodyfont\color{darktext}#1}\hfill{\fontsize{9.5pt}{12.5pt}\bodyfont\color{darktext}#2}\par\nopagebreak}

\newcommand{\skillline}[2]{\par\noindent\hangindent=1.2em{\fontsize{10pt}{12.6pt}\bodyfont\color{text}\textbf{#1:} #2}\par\vspace{1.5pt}}
\newcommand{\pubitem}[1]{\par\noindent\hangindent=1.2em{\fontsize{9.5pt}{12pt}\bodyfont\color{text}#1}\par\vspace{2pt}}
\newcommand{\eduline}[3]{\par\noindent{\fontsize{10pt}{12.6pt}\bodyfont\color{text}\textbf{#1}, #2}\hfill{\fontsize{9.5pt}{12.6pt}\bodyfont\color{darktext}#3}\par\vspace{1pt}}
\newcommand{\me}{\textbf{C.~E.~Tejada}}

% ---------------------------------------------------------------------
% 6. Cover-letter paragraph bank: \letterpara{id}{role types}{text}, \usepara{id}
% ---------------------------------------------------------------------
\def\paralist{}
\newcommand{\letterpara}[3]{\csgdef{para@#1}{#3}\csgdef{paratags@#1}{#2}\listgadd\paralist{#1}}
\newcommand{\usepara}[1]{\ifcsdef{para@#1}{\csuse{para@#1}}{\slot{missing paragraph "#1"}}}
